\documentclass[10pt,letterpaper]{article}

\usepackage[margin=0.48in]{geometry}
\usepackage[T1]{fontenc}
\usepackage{lmodern}
\usepackage{xcolor}
\usepackage{enumitem}
\usepackage{tabularx}
\usepackage[hidelinks]{hyperref}
\usepackage{titlesec}
\usepackage{microtype}

\definecolor{accent}{HTML}{2F6B4F}
\setlength{\parindent}{0pt}
\setlength{\parskip}{0pt}
\setlist[itemize]{leftmargin=1.15em,itemsep=0.7pt,topsep=1pt,parsep=0pt}
\titleformat{\section}{\large\bfseries\color{accent}}{}{0pt}{}[\vspace{-2pt}\titlerule]
\titlespacing*{\section}{0pt}{4pt}{2pt}
\pagestyle{empty}

\newcommand{\role}[4]{%
  \textbf{#1}\hfill #2\\
  \textit{#3}\hfill \textit{#4}\vspace{-3pt}
}

\begin{document}
\fontsize{9}{10.4}\selectfont

\begin{center}
  {\Large\bfseries Piter Z. Garcia Bautista}\\[1pt]
  {\large Program Evaluation, Inclusive Education, Disability, and Data Systems}\\[2pt]
  Rochester, NY \enspace|\enspace Available to relocate after December 2026\\
  \href{mailto:pzg8794@rit.edu}{pzg8794@rit.edu}
  \enspace|\enspace +1 (646) 288-3300
  \enspace|\enspace \href{https://linkedin.com/in/garciapiterz}{linkedin.com/in/garciapiterz}\\
  \href{https://github.com/pzg8794}{github.com/pzg8794}
  \enspace|\enspace \href{https://garciapiterz.com}{garciapiterz.com}
\end{center}

\section*{Profile}
Multilingual graduate educator and data practitioner connecting program
evaluation, inclusive computer science, disability access, responsible AI,
and clear cross-functional communication. Experienced in designing data
workflows, validating results, documenting requirements, translating
technical findings, and developing accessible learning for varied
participants. Available for the Carman International Fellowship's one-year
placement beginning in January 2027.

\section*{Relevant Experience}
\role{Computer Science Teacher Candidate}{2025--Present}
{University of Rochester / Pine Brook Elementary School}{Rochester, NY}
\begin{itemize}
  \item Co-design standards-aligned learning in coding, robotics, digital
    fluency, and responsible AI for learners with varied strengths and needs.
  \item Use Universal Design for Learning and multiple participation routes to
    expand access while preserving rigorous goals and respectful expectations.
\end{itemize}

\role{Graduate Assistant, Quantum and AI Research}{Fall 2025--Present}
{Rochester Institute of Technology, Software Engineering}{Rochester, NY}
\begin{itemize}
  \item Build Python workflows for data collection, automated validation,
    structured logging, reproducible evaluation, and comparison-ready
    reporting.
  \item Translate complex requirements and large-scale findings into clear
    analyses, presentations, documentation, and manuscript artifacts.
\end{itemize}

\role{Data Solutions Engineer}{Sep 2022--Aug 2024}
{VEDADATA Inc.}{New York, NY}
\begin{itemize}
  \item Designed maintainable data-ingestion, preprocessing, validation, and
    analytics workflows for production and business applications.
  \item Collaborated across technical and nontechnical teams to clarify needs,
    document findings, and improve reliability and usability.
\end{itemize}

\role{AI Data Solutions Engineer}{Jan 2021--Sep 2022}
{VIOME Inc.}{New York, NY}
\begin{itemize}
  \item Developed healthcare-oriented machine-learning workflows for
    preprocessing, feature engineering, evaluation, and accessible reporting.
\end{itemize}

\section*{Education}
\textbf{M.S., Teaching Computer Science K--12}, University of Rochester
\hfill Expected December 2026\\
GPA: 3.9+; NSF Robert Noyce Scholar; Advanced Certificate in Leadership in
Disability and Inclusive Practices

\textbf{M.S., Data Science}, Rochester Institute of Technology
\hfill Expected December 2026\\
GPA: 3.9+; applied AI, bioinformatics, neural networks, and reproducible data systems

\textbf{M.S., Computer Science}, Rochester Institute of Technology
\hfill 2015\\
Artificial intelligence, big data analytics, and computer graphics

\textbf{B.S., Computer Engineering Technology}, Farmingdale State College
\hfill 2012

\section*{Selected Program and Access Work}
\textbf{EQUITAS and the Puzzle Plan.}
Use EQUITAS, the equity-focused lens within the interdisciplinary Puzzle Plan,
to connect responsible data systems, disability justice, communication access,
and inclusive education.

\textbf{Program Monitoring and Reproducible Evaluation.}
Create structured datasets, validation checks, traceable experiment logs, and
clear comparison artifacts that help collaborators identify gaps and make
evidence-informed decisions.

\textbf{Accessible Technical Communication.}
Develop high-contrast, multimodal, self-contained learning and reporting
materials that support different participation routes without ranking one
communication mode over another.

\section*{Skills and Languages}
\begin{tabularx}{\textwidth}{@{}>{\bfseries}l X@{}}
Program support & Monitoring and evaluation, data collection, quality checks, documentation, needs clarification, action tracking, report and presentation development\\
Education and access & Inclusive CS curriculum, UDL, differentiated participation, disability leadership, assessment design\\
Technical & Python, SQL, R, Java, data analysis, machine learning, dashboards and visualization, Git/GitHub, reproducible workflows\\
Languages & Spanish: native/fluent; English: fluent\\
\end{tabularx}

\section*{Recognition}
NSF Robert Noyce Scholar (2024--2026)
\enspace|\enspace IEEE Alumni Member (2009--Present)
\enspace|\enspace RIT/MESCYT merit-based graduate funding

\end{document}
